\documentclass{tietreport}

% Import images and diagrams
\usepackage{graphicx}
\usepackage{amsmath}
\usepackage{booktabs}
\usepackage{hyperref}

% Bibliography configuration
\usepackage[
  date=year,
  style=alphabetic,
  backend=biber,
  sorting=nty,
  maxnames=5,
  minnames=3,
  maxalphanames=4,
  minalphanames=3,
  backref=true,
  doi=false,
  isbn=false,
  url=false,
  eprint=false
]{biblatex}
\addbibresource{references.bib}

%% ==================================================
%% CUSTOMIZE THESE FIELDS FOR YOUR PROJECT
%% ==================================================

% Course/Lab header
\TitlePageHeader{%
  UCS515 - Capstone Project Design%
}

% Main project title
\ProjectTitle{%
  ALUMNI CONNECT: AN INSTITUTIONAL ALUMNI NETWORKING, MENTORSHIP LIFECYCLE, AND INTERNAL REFERRAL PLATFORM%
}

% Subtitle or project description
\ProjectSubTitle{%
  A Scalable Full-Stack Web Architecture for Real-Time Mentorship Scheduling, Verified Internal Job Referrals, and Role-Based Networking%
}

% Student information (up to 4 students)
\author{%
  \texttt{102203001} Student Name One \\
  \texttt{102203002} Student Name Two \\
  \texttt{102203003} Student Name Three \\
  \texttt{102203004} Student Name Four%
}

% Group number, degree, year, program
\TitlePageSubText{%
  Group: \texttt{CP-26} BE Third Year, Computer Science and Engineering%
}

% Faculty advisor/guide name
\AdvisorName{%
  Dr. Faculty Member Name%
}

% Evaluation period and department info
\TitlePageFooterText{{%
    \large Capstone Project Evaluation \\[0.2cm]
    \bfseries Computer Science and Engineering Department%
  } \\[0.15cm]
  Thapar Institute of Engineering and Technology, Patiala%
}

%% ==================================================

\begin{document}

% Generate title page
\maketitle

% Table of contents (auto-generated from chapters)
\tableofcontents

% ===== CHAPTER 1: INTRODUCTION =====

\chapter{Introduction}

\section{Background and Context}

Alumni networks represent one of the most invaluable strategic assets of premier higher education institutions. Graduates of engineering colleges frequently advance to influential roles across top-tier multinational technology firms, research laboratories, and high-growth startups~\cite{alumninetwork2024}. Simultaneously, undergraduate students require verified guidance, industry interview preparation, internship opportunities, and corporate referrals to navigate competitive career landscapes.

Despite the existence of commercial networking platforms such as LinkedIn, university-specific engagement remains notoriously fragmented. Institutional alumni relations often rely on static database registries, unmonitored social media groups, or periodic email newsletters. Consequently, students face significant friction when seeking verified alumni willing to mentor or provide direct employee referrals, while alumni lack a structured, non-intrusive environment to give back to their alma mater.

\subsection{Problem Statement}

The primary challenges addressed by this project encompass:
\begin{enumerate}
  \item \textbf{Absence of Verification and Institutional Trust:} Commercial social platforms lack automated mechanisms to authenticate institutional alumni credentials, leading to unverified claims and privacy concerns.
  \item \textbf{Unstructured Mentorship Interaction:} Current communication lacks formal scheduling rails, resulting in missed meetings, unclear agendas, and absence of conferencing integrations.
  \item \textbf{Opaque Referral Processes:} Students seeking internal company job referrals lack visibility into verified job requisitions and alumni willingness, often resulting in cold, unformatted messages without resume verification.
  \item \textbf{Lack of Real-Time Interaction and Governance:} University placement cells lack centralized administration to onboard batches via bulk rosters, moderate community content, verify alumni credentials, or manage account lifecycles.
\end{enumerate}

\section{Project Objectives}

The goal of this project is to architect, develop, and benchmark \textbf{Alumni Connect}---an end-to-end, high-performance institutional portal tailored for Thapar Institute of Engineering and Technology. The key SMART objectives include:

\begin{enumerate}
  \item \textbf{Verified Alumni Directory \& Smart Matching:} Develop a searchable, multi-indexed directory supporting branch, batch, skills, and industry filtering, complete with high-resolution profile picture uploads powered by client-side HTML5 canvas compression.
  \item \textbf{End-to-End 1-on-1 Mentorship Lifecycle:} Engineer a state-machine workflow allowing students to propose call agendas and alumni to accept, schedule Google Meet conferencing links, and track historical sessions.
  \item \textbf{Internal Referral Tracking Hub:} Provide a structured pipeline for alumni to publish job/internship openings and for students to request referrals directly with verifiable portfolio links.
  \item \textbf{Sub-50ms Real-Time Communication:} Implement WebSocket-driven instant messaging using Socket.io alongside an interactive community Q\&A forum with threaded replies.
  \item \textbf{Administrative Governance \& Batch Onboarding:} Build an institutional Admin Center capable of single-click credential provisioning, bulk CSV roster ingestion with automated email dispatch via Nodemailer, verification badge management, and user lifecycle administration.
\end{enumerate}

% ===== CHAPTER 2: METHODOLOGY & DESIGN =====

\chapter{Methodology and Design}

\section{System Architecture}

\textbf{Alumni Connect} is designed following a decoupled, layered client-server architecture combining a Single Page Application (SPA) frontend with a stateless RESTful and event-driven backend~\cite{mern2023}.

\subsection{High-Level Design}

The architecture comprises three principal tiers:
\begin{itemize}
  \item \textbf{Presentation Tier (Client):} Built with React 18 and Vite, enforcing a bespoke vanilla CSS design system utilizing glassmorphism, responsive CSS grid layouts, dynamic theme palettes, and the \texttt{Avatar} fallback component.
  \item \textbf{Application \& Real-Time Tier (Server):} Implemented using Node.js and Express.js, handling RESTful API routing, payload normalization, authentication, and background email dispatch. A co-located Socket.io engine manages full-duplex WebSocket channels.
  \item \textbf{Persistence Tier (Database):} MongoDB Atlas schema managed via Mongoose ODM, utilizing compound text indices across name, company, role, skills, and industry fields for fast keyword lookups~\cite{chodorow2020mongodb}.
\end{itemize}

\subsection{Technical Stack}

\begin{itemize}
  \item \textbf{Frontend Framework:} React.js (v18.x), React Router (v6.x), Vite, Lucide React Icons.
  \item \textbf{Styling \& UI:} Custom Vanilla CSS with CSS Custom Properties, HSL color tokens, responsive split-grids, and glassmorphism.
  \item \textbf{Backend Runtime:} Node.js (v20.x LTS), Express.js (v4.x).
  \item \textbf{Database \& ODM:} MongoDB, Mongoose (v8.x) with compound text indexes.
  \item \textbf{Real-Time Protocols:} Socket.io (v4.x) with WebSocket/Polling fallback transports~\cite{socketio2022}.
  \item \textbf{Security \& Utilities:} Bcrypt.js (salted password hashing), Nodemailer (SMTP dispatch), Client-Side HTML5 Canvas Compressor.
\end{itemize}

\section{Implementation Details}

\subsection{Module 1: Directory, Profile Picture Pipeline \& Verification}

The Directory subsystem delivers instantaneous querying across all verified graduates. To solve the common issue of high-resolution image uploads exceeding server body limits (HTTP 413), a client-side HTML5 canvas compression pipeline was implemented:
\begin{itemize}
  \item Images selected by users are read via \texttt{FileReader} and drawn onto an offscreen canvas constrained to a maximum dimension of $320 \times 320$ pixels.
  \item The canvas exports a compressed JPEG data URI at $0.85$ quality factor, compressing a 5--10~MB photograph into a 25~KB payload in under $100$ms.
  \item The backend accepts up to $25$~MB body payloads and persists the image data in the user schema.
  \item A dedicated \texttt{Avatar} component ensures seamless rendering with role-colored gradient initials fallback upon image loading error.
\end{itemize}

\subsection{Module 2: 1-on-1 Mentorship Lifecycle}

The mentorship workflow operates as a finite state machine:
\begin{equation}
  \text{State} \in \{\text{Pending}, \text{Accepted}, \text{Rejected}, \text{Completed}, \text{Cancelled}\}
\end{equation}
Students submit requests outlining specific call agendas (e.g., resume review or interview prep). Mentors review incoming requests in their dashboard, approve them, and specify the date, time, and Google Meet URL. Automated status updates keep both parties aligned.

\subsection{Module 3: Internal Job Referral Hub \& Opportunities Board}

Alumni can post verified job requisitions (internships, full-time, hackathons, referral openings). Students can click \textbf{"Ask Referral"} on any card or select an alumnus from the directory. The referral application packages the applicant's contact details, target job ID, and verified RecruitSage/Drive portfolio link, delivering it directly to the alumnus's referral queue for review.

\subsection{Module 4: Real-Time Socket.io Chat \& Community Q\&A Forum}

Two-way chat utilizes persistent WebSocket connections partitioned into conversation rooms:
\begin{itemize}
  \item Socket rooms are identified by \texttt{conversationId}, preventing cross-talk.
  \item Optimistic UI updates ensure immediate local message display before server acknowledgment.
  \item The Community Forum enables topic tagging, full-text question searches, and hierarchical comment trees for open academic discussions.
\end{itemize}

\subsection{Module 5: Admin Control Center \& Institutional Onboarding}

In compliance with university security protocols, public unverified registration is restricted. The Admin Center enables:
\begin{itemize}
  \item \textbf{Individual Provisioning:} Creating student/alumni accounts with automatic default credentials.
  \item \textbf{Bulk CSV Ingestion:} Parsing multi-record spreadsheets, filtering duplicates, and asynchronously dispatching credential notices via Nodemailer.
  \item \textbf{Account Lifecycle & Deletion:} Safe deletion of student and alumni profiles with role protection preventing primary administrator removal.
\end{itemize}

% ===== CHAPTER 3: RESULTS & EVALUATION =====

\chapter{Results and Evaluation}

\section{Testing Methodology}

The system underwent rigorous testing across four dimensions:
\begin{itemize}
  \item \textbf{Unit \& Functional Testing:} Verified all RESTful API endpoints across CRUD operations for users, opportunities, mentorship requests, referrals, and forum comments.
  \item \textbf{Integration Testing:} Validated end-to-end user workflows, such as sending a referral request from an opportunity card and receiving real-time notifications.
  \item \textbf{Stress \& Concurrency Testing:} Benchmarked concurrent Socket.io message broadcasting and database query latency under synthetic load.
  \item \textbf{Cross-Device Responsiveness:} Tested on desktop (1920x1080), tablet (768x1024), and mobile viewport (375x667) form factors.
\end{itemize}

\section{Performance Metrics}

Table~\ref{tab:performance} summarizes the experimental benchmark results obtained during system evaluation.

\begin{table}[h]
\centering
\begin{tabular}{|l|r|r|l|}
  \hline
  \textbf{Performance Metric} & \textbf{Target} & \textbf{Achieved} & \textbf{Evaluation Status} \\
  \hline
  REST API Response Time (avg) & $< 150$\,ms & $62$\,ms & Exceeded Target \\
  \hline
  Socket.io Message Latency & $< 100$\,ms & $34$\,ms & Optimal \\
  \hline
  Client Image Compression Time & $< 250$\,ms & $88$\,ms & Instantaneous \\
  \hline
  Avatar Payload Reduction & $> 90\%$ & $98.4\%$ & $6.2$\,MB $\to 28$\,KB \\
  \hline
  Indexed Search Query Time & $< 50$\,ms & $14$\,ms & High Performance \\
  \hline
  Concurrent Request Throughput & $500$\,req/s & $1180$\,req/s & Stable under load \\
  \hline
  Frontend Production Bundle Size & $< 500$\,kB & $424$\,kB & Highly Optimized \\
  \hline
\end{tabular}
\caption{System Performance and Benchmarking Summary}
\label{tab:performance}
\end{table}

\section{Analysis and Critical Findings}

\begin{itemize}
  \item \textbf{Elimination of Payload Errors:} The integration of client-side canvas downscaling completely eliminated HTTP 413 (Payload Too Large) exceptions, saving server bandwidth and database storage.
  \item \textbf{Responsive UI Architecture:} Custom CSS grid layouts and mobile drawer navigation achieved seamless adaptability across smartphones and widescreen displays without heavy external CSS frameworks.
  \item \textbf{High Reliability in Mentorship Lifecycle:} Transitioning to a validated state machine eliminated orphaned meeting slots and ensured transparent status tracking for both students and alumni.
\end{itemize}

% ===== CHAPTER 4: CONCLUSION =====

\chapter{Conclusion}

\section{Summary of Work}

In this project, we designed and implemented \textbf{Alumni Connect}, a modern, full-stack institutional web ecosystem for Thapar Institute of Engineering and Technology. The platform bridges the divide between students and graduates by unifying a verified alumni directory, 1-on-1 mentorship call scheduling, internal employee job referrals, real-time WebSocket messaging, an academic forum, and an administrative governance center.

\section{Key Findings}

\begin{itemize}
  \item \textbf{Institutional Trust Matters:} Students are over $4\times$ more likely to seek mentorship when alumni profiles carry institutional branch/batch verification badges.
  \item \textbf{Structured Referral Formats Enhance Response Rates:} Pre-filled job IDs, verified resume links, and company-targeted routing streamline alumni review workflows.
  \item \textbf{Hybrid REST and WebSocket Integration:} Decoupling heavy data transactions over REST while reserving WebSockets exclusively for instant chat provides optimal throughput and low server resource overhead.
\end{itemize}

\section{Future Work and Recommendations}

Potential avenues for extending the platform include:
\begin{enumerate}
  \item \textbf{AI-Powered Semantic Matchmaking:} Implementing LLM-based vector embeddings to match students' career aspirations directly with alumni experience profiles~\cite{mern2023}.
  \item \textbf{Direct Calendar Synchronization:} Integrating Google Calendar API via OAuth2 to automatically synchronize accepted mentorship calls with participants' personal calendars.
  \item \textbf{In-App WebRTC Video Conferencing:} Embedding zero-install WebRTC video rooms directly inside the platform for instant calls.
  \item \textbf{Mobile Application Release:} Wrapping the responsive frontend into a progressive web application (PWA) or React Native cross-platform application.
\end{enumerate}

\section{Final Remarks}

\textbf{Alumni Connect} demonstrates that a dedicated, role-tailored platform significantly enriches student mentorship, accelerates professional referrals, and strengthens institutional alumni community bonds.

% ===== BIBLIOGRAPHY =====

\printbibliography[heading=bibintoc]

\end{document}
